A recovery curve can measure sampling performance while reproducing a specified
empirical reference. The residual makes that equality visible, but neither
equality nor departure ranks estimator quality. Our patient simulations show
that the comparison depends on the reference, treatment-assignment design and
which information censoring hides. The clean control makes the limitation
explicit: identical point estimates can have very different interval coverage,
and a legitimate estimator can depart from the declared reference.

The motivating recovery curve did not justify the cell-count rule: direct
coverage and discrimination evaluation returned no qualifying cutpoint or
configuration-dependent candidates. Reusing identical draws across count grids
confirmed that both available counts and aggregation affect those candidates.
The smaller cohort's candidate changed with bootstrap budget and, separately,
patient omission. The synthetic benchmark
(Appendix~\ref{app:bench}) likewise found that an existence check prevented the
same false interpretations with fewer refusals under the tested Poisson model.
Pre-specify the reference, report exclusions and residual scale, and investigate
numerical equality algebraically. Evaluate requested-parameter error, interval
coverage and the decision rule separately.
